% ======================================================================
% Archivo: estado_del_arte.tex
% Capítulo: Estado del Arte (reestructurado)
% Capítulo: Marco Teórico (esqueleto preservado)
% ======================================================================

% ======================================================================
% ESTADO DEL ARTE
% ======================================================================
\section{Estado del Arte}
\label{ch:estado-del-arte}

El presente capítulo revisa la literatura científica relacionada con la sustentación técnica de esta tesina: la representación factorizada de MDPs y el problema de las variables de estado, el uso de la Programación Lógica Probabilística como lenguaje de especificación para la toma de decisiones secuenciales, y los mecanismos de exclusión mutua y compilación de conocimiento que permiten la inferencia eficiente. El análisis de estos trabajos permite trazar la trayectoria que conduce al vacío de investigación que la presente tesina aborda: la ausencia de soporte nativo para fluentes de estado multivaluados en frameworks de resolución exacta de MDPs basados en PLP.

% ------------------------------------------------------------------
\subsection{MDPs Factorizados y el Problema de la Representación de Estados}
\label{sec:ea-factorizados}

En la formulación clásica de un MDP, el espacio de estados se trata como un conjunto plano donde cada estado es una entidad atómica e indivisible \cite{puterman2014markov}. Esta representación se vuelve poco práctica conforme aumenta el número de variables descritas: si un dominio involucra $n$ variables binarias, el número total de estados crece como $2^n$, lo que Bellman denominó la \textit{maldición de la dimensionalidad} \cite{bellman1957dynamic}.

Para mitigar este crecimiento exponencial, Boutilier, Dearden y Goldszmidt propusieron la representación factorizada de MDPs, donde las transiciones se codifican mediante Redes Bayesianas Dinámicas (DBNs) y las funciones de valor se manipulan como árboles de decisión \cite{boutilier2000stochastic}. En esta representación, cada variable de estado posee su propia tabla de probabilidad condicional, y los algoritmos de programación dinámica operan sobre agregaciones de estados en lugar de iterar sobre cada estado individual. Su algoritmo \textit{Structured Policy Iteration} (SPI) demostró que era posible resolver MDPs sin enumerar explícitamente el espacio de estados, siempre que la estructura del problema exhibiera independencias condicionales aprovechables.

En la misma línea, Hoey \textit{et al.} desarrollaron SPUDD, un algoritmo de iteración de valor que reemplaza los árboles de decisión por Diagramas de Decisión Algebraicos (ADDs) \cite{hoey1999spudd}. Los ADDs, al ser grafos dirigidos acíclicos en lugar de árboles, permiten compartir subgrafos isomorfos, logrando reducciones de hasta 30 veces en el número de nodos respecto a las representaciones arbóreas. SPUDD demostró la resolución exacta de MDPs con hasta 63 millones de estados.

Tanto SPI como SPUDD comparten una restricción, ya que operan exclusivamente sobre variables booleanas. Cuando el dominio contiene variables con más de dos valores, el modelador debe descomponer cada variable multivaluada en múltiples variables binarias. Hoey \textit{et al.} reconocen explícitamente esta limitación al señalar que variables ternarias deben expandirse en dos binarias, lo cual agranda el espacio de estados por un factor de $4/3$ por cada variable expandida e introduce estados semánticamente imposibles \cite{hoey1999spudd}. Este problema de la binarización forzada es precisamente la limitación que la presente tesina aborda en el contexto de MDP-ProbLog.

En el ámbito de los lenguajes de descripción para planificación, RDDL (\textit{Relational Dynamic Influence Diagram Language}) \cite{sanner2010rddl} sí contempla fluentes multivaluados de manera nativa, incluyendo variables enumeradas, enteras y continuas. No obstante, RDDL está fundamentado en redes bayesianas dinámicas proposicionales, lo que impide la representación de dependencias cíclicas que surgen naturalmente en dominios relacionales. MDP-ProbLog, en contraste, hereda de la programación lógica la capacidad de expresar este tipo de relaciones, pero carece del soporte multivaluado que RDDL ofrece.

% ------------------------------------------------------------------
\subsection{Programación Lógica Probabilística para la Resolución de MDPs}
\label{sec:ea-plp-mdps}

La Programación Lógica Probabilística (PLP) combina la capacidad de representación relacional de la lógica de primer orden con el manejo formal de la incertidumbre mediante la teoría de la probabilidad. De Raedt y Kimmig establecieron los fundamentos unificadores de este paradigma, demostrando cómo la \textit{semántica de distribución} permite definir una distribución de probabilidad sobre los mundos posibles de un programa lógico \cite{deraedt2015probabilistic}. Bajo esta semántica, cada hecho probabilístico actúa como una variable aleatoria independiente cuya realización determina un programa lógico determinista; la probabilidad de una consulta es entonces la suma de las probabilidades de todos los mundos en los que dicha consulta es verdadera.

ProbLog, el lenguaje que implementa esta semántica, transforma el problema de inferencia probabilística en un problema de Conteo de Modelos Ponderados (WMC) sobre fórmulas booleanas \cite{fierens2015inference, deraedt2007problog}. Esto permite aprovechar técnicas de compilación de conocimiento para lograr inferencia exacta en tiempo polinomial respecto al tamaño del circuito compilado.

La extensión de ProbLog hacia la toma de decisiones se materializó primero en DTProbLog \cite{vandenbroeck2010dtproblog}, que añadió al lenguaje la capacidad de modelar problemas de decisión episódicos: dado un conjunto de variables de decisión y atributos de utilidad, el sistema encuentra la estrategia que maximiza la utilidad esperada. Sin embargo, DTProbLog está diseñado para problemas de un solo paso y no soporta la resolución de MDPs de horizonte infinito, donde es necesario optimizar una política completa a lo largo del tiempo.

Como derivado de este trabajo, Bueno \textit{et al.} propusieron MDP-ProbLog, el cual extiende la sintaxis de ProbLog con predicados reservados para declarar fluentes de estado, acciones y utilidades, y emplea programas de distribución de transición de dos pasos temporales para representar la dinámica del MDP \cite{bueno2016mdp}. El \textit{solver} implementa iteración de valor sobre el circuito compilado mediante WMC, demostrando formalmente la equivalencia entre ejecutar el programa probabilístico y realizar un \textit{backup} de Bellman. Una ventaja distintiva de MDP-ProbLog frente a lenguajes como RDDL es su capacidad de representar dependencias cíclicas entre predicados aterrizados gracias a la resolución SLD de Prolog \cite{bueno2016mdp}.

% ------------------------------------------------------------------
\subsection{Disyunciones Anotadas y Compilación de Conocimiento}
\label{sec:ea-ads-compilacion}

El mecanismo que esta tesina propone para resolver la restricción binaria de MDP-ProbLog se fundamenta en las Disyunciones Anotadas (ADs). En la literatura de representación del conocimiento, Vennekens, Verbaeten y Bruynooghe introdujeron los Logic Programs with Annotated Disjunctions (LPADs) como una solución elegante para modelar eventos probabilísticos con múltiples resultados mutuamente excluyentes. El principal resultado de su trabajo demostró que es posible codificar la exclusión mutua de manera nativa en la cabeza de las reglas lógicas, eliminando la necesidad de escribir restricciones manuales complejas y mitigando la generación de estados semánticamente imposibles.

Adicionalmente, Vennekens et al. demostraron empírica y formalmente que las ADs admiten una interpretación causal natural para problemas de toma de decisiones. Esta característica es crucial, ya que permite mapear de forma directa la semántica del lenguaje a las transiciones estocásticas de un MDP: dado un estado y una acción (la causa), el sistema transita hacia uno de múltiples estados posibles (los efectos). En esta misma línea, demostraron la equivalencia expresiva entre su propuesta y el Independent Choice Logic (ICL) de Poole, consolidando a las ADs como un estándar más natural para representar efectos indeterminados.

Ahora bien, para que la inferencia probabilística sobre programas con ADs sea computacionalmente tratable en la práctica, sistemas como ProbLog emplean un flujo de compilación de conocimiento. El programa lógico aterrizado se transforma primero en una fórmula proposicional ponderada, y posteriormente se compila en un circuito lógico. En este ámbito, Darwiche y Marquis establecieron la taxonomía de los lenguajes de compilación, identificando que estructuras como las d-DNNFs (Deterministic Decomposable Negation Normal Forms) permiten realizar inferencia exacta en tiempo lineal respecto al tamaño del circuito compilado \cite{darwiche2002knowledge}. Posteriormente, el mismo autor introdujo los Diagramas de Decisión Sentencial (SDDs), una representación canónica que generaliza los clásicos BDDs y resulta estrictamente más sucinta \cite{darwiche2011sdd}. En la práctica, el motor subyacente de ProbLog2, brinda la flexibilidad de compilar el modelo hacia cualquiera de estas dos estructuras \cite{vlasselaer2016problog2}. Es precisamente sobre estos circuitos donde el sistema ejecuta el Conteo de Modelos Ponderados (WMC) para calcular las probabilidades de transición. Por lo tanto, cualquier extensión a la capa de representación, como la integración de fluentes multivaluados propuesta en este trabajo, debe garantizar su compatibilidad con esta forma de compilación para preservar la eficiencia computacional.

\clearpage

% ------------------------------------------------------------------
\subsection{Análisis Comparativo}
\label{sec:ea-comparativa}

La Tabla~\ref{tab:estado_arte_comparativa} presenta una comparativa de los principales enfoques revisados en este capítulo. Se analizan diez trabajos representativos agrupados por familia metodológica. Para cada uno se evalúan cuatro dimensiones: el tipo de variables de estado que soporta, el mecanismo de inferencia empleado, la capacidad del lenguaje para expresar relaciones entre entidades, y la forma en que se garantiza la exclusión mutua entre los valores de una misma variable.

\begin{table}[H]
    \centering

    \caption{Análisis comparativo del estado del arte. Se evalúan cuatro dimensiones para cada trabajo: el tipo de variables de estado soportadas, el mecanismo de inferencia, la capacidad relacional del lenguaje, y la garantía de exclusión mutua entre valores de una misma variable.}
    \label{tab:estado_arte_comparativa}
    
    \footnotesize 
    \setlength{\tabcolsep}{4pt} 
    \renewcommand{\arraystretch}{1.2} 
    
    \begin{tabularx}{\textwidth}{@{} 
        >{\raggedright\arraybackslash}p{3cm} 
        >{\raggedright\arraybackslash}X 
        >{\raggedright\arraybackslash}X 
        >{\raggedright\arraybackslash}X 
        >{\raggedright\arraybackslash}p{2.5cm} @{}} 
        \toprule
        \textbf{Trabajo} &
        \textbf{Variables de estado} &
        \textbf{Inferencia} &
        \textbf{Expresividad relacional} &
        \textbf{Exclusión mutua} \\
        \midrule
        
        \multicolumn{5}{@{}l}{\textbf{\textit{MDPs factorizados proposicionales}}} \\
        Boutilier \textit{et al.} \cite{boutilier2000stochastic} & Booleanas & Exacta (árboles) & Proposicional (DBNs) & Manual \\
        Boutilier, Dean y Hanks \cite{boutilier1999decision} & Booleanas & Exacta (prog. dinámica) & Proposicional & Manual \\
        Hoey \textit{et al.} \cite{hoey1999spudd} & Booleanas (binarización) & Exacta (ADDs) & Proposicional (DBNs) & Manual \\
        \addlinespace
        
        \multicolumn{5}{@{}l}{\textbf{\textit{Lenguajes de planificación}}} \\
        Sanner \cite{sanner2010rddl} & Multival., continuas & Exacta / Aprox. & Relacional (DBNs) & Nativa por tipo \\
        \addlinespace
        
        \multicolumn{5}{@{}l}{\textbf{\textit{Programación Lógica Probabilística}}} \\
        De Raedt y Kimmig \cite{deraedt2015probabilistic} & Booleanas & Exacta (WMC) & Relacional (ProbLog) & No aplica \\
        Fierens \textit{et al.} \cite{fierens2015inference} & Booleanas & Exacta (WMC) & Relacional (ProbLog) & No aplica \\
        Nitti \textit{et al.} \cite{nitti2016planning} & Discretas, continuas & Aprox. (muestreo) & Relacional (DDC) & Implícita \\
        Riguzzi y Swift \cite{riguzzi2013well} & Booleanas & Exacta (tabling) & Relacional & No aplica \\
        \addlinespace
        
        \multicolumn{5}{@{}l}{\textbf{\textit{Exclusión mutua y semántica causal}}} \\
        Vennekens \textit{et al.} \cite{vennekens2009cplogic} & Multivaluadas (ADs) & Exacta & Relacional & Nativa  \\
        Poole \cite{poole1997independent} & Booleanas (ICL) & Exacta & Relacional & Nativa (elecciones) \\
        \addlinespace
        
        \multicolumn{5}{@{}l}{\textbf{\textit{MDP-ProbLog}}} \\
        Bueno \textit{et al.} \cite{bueno2016mdp} & Booleanas & Exacta (WMC/SDDs) & Relacional (ProbLog) & Manual \\
        \textbf{Propuesta} & \textbf{Multivaluadas} & \textbf{Exacta (WMC/SDDs)} & \textbf{Relacional (ProbLog)} & \textbf{Nativa (ADs)} \\
        \bottomrule
    \end{tabularx}
\end{table}

Como se observa en la tabla, ninguno de los enfoques existentes integra simultáneamente las cuatro capacidades que la propuesta de esta tesina reúne: variables de estado multivaluadas con exclusión mutua nativa, inferencia exacta mediante WMC sobre circuitos compilados, y expresividad relacional completa heredada de la programación lógica.

% ------------------------------------------------------------------
\subsection{Conclusión del Estado del Arte}
\label{sec:ea-conclusion}

La revisión de la literatura permite identificar tres hallazgos principales. En primer lugar, la representación factorizada de MDPs ha demostrado ser una estrategia efectiva para mitigar la maldición de la dimensionalidad, permitiendo resolver problemas con millones de estados sin enumeración explícita \cite{boutilier2000stochastic, hoey1999spudd}. Sin embargo, los enfoques proposicionales más consolidados operan exclusivamente sobre variables booleanas, y cuando el dominio requiere variables multivaluadas, la binarización forzada introduce estados semánticamente imposibles y aumenta artificialmente el tamaño de los circuitos compilados.

En segundo lugar, la Programación Lógica Probabilística ofrece ventajas significativas sobre los formalismos proposicionales al permitir representar relaciones simétricas, transitivas y cíclicas entre las variables del dominio \cite{bueno2016mdp, deraedt2015probabilistic}. El framework MDP-ProbLog materializó esta ventaja al combinar la expresividad de ProbLog con la resolución exacta de MDPs de horizonte infinito mediante WMC \cite{fierens2015inference}. No obstante, su arquitectura restringe los fluentes de estado a variables de Bernoulli independientes, lo que limita la naturalidad del proceso de modelado y delega en el usuario la responsabilidad de codificar manualmente las restricciones de exclusión mutua.

En tercer lugar, las Disyunciones Anotadas constituyen un mecanismo formal que garantiza la exclusión mutua de eventos probabilísticos a nivel semántico \cite{vennekens2009cplogic, poole1997independent}. Este constructo ya forma parte del lenguaje ProbLog y es utilizado internamente por su motor de inferencia \cite{vlasselaer2016problog2}, pero su integración en la capa de representación de fluentes de estado dentro de MDP-ProbLog no ha sido explorada ni implementada.

En consecuencia, la presente tesina se posiciona en la intersección de estos tres ejes al proponer la extensión de MDP-ProbLog para soportar fluentes de estado multivaluados mediante Disyunciones Anotadas. Esta integración busca preservar las garantías de inferencia exacta del sistema, eliminar la necesidad de binarización manual, y habilitar una representación del espacio de estados que sea tanto más compacta como semánticamente íntegra. El siguiente capítulo establece los fundamentos formales --- desde la programación lógica hasta la compilación de conocimiento --- que el lector necesitará para comprender el diseño de la extensión propuesta.
